We study a narrow but practical question: if a user exports music listening history and asks a frontier language model for recommendations, can prompting alone match or beat lightweight Spotify-style recommenders? We evaluate this question as closed-set ranking on two public tasks: \lastfm next-track prediction and \mpd playlist continuation. Each test example contains one held-out relevant track and 19 negatives, so classical baselines and language models solve the same 20-way ranking problem. We compare two prompt-only LLMs, \gptfourone and \claudesonnet, against two classical baselines: a popularity ranker and a personalized cooccurrence heuristic.

The strongest results come from the simplest prompt. On \lastfm, \gptfourone with the \minimalprompt prompt reaches 0.589 \mrr and 0.621 \ndcg, compared with 0.282 \mrr and 0.246 \ndcg for the strongest classical baseline. On \mpd, the same setup reaches 0.329 \mrr and 0.379 \ndcg, compared with 0.094 \mrr and 0.061 \ndcg for cooccurrence. These gains are statistically significant under paired Wilcoxon signed-rank tests with Benjamini-Hochberg correction. Adding compact profile hints does not help and is reliably worse on \mpd.

Our main conclusion is narrow but clear: prompt-only frontier LLMs can outperform lightweight history-based music recommendation baselines on public offline ranking tasks. At the same time, this result does not establish that prompting beats Spotify's production recommender, because we evaluate public datasets, fixed candidate pools, and simpler baselines than a production system is likely to use.
